\documentclass[12pt,a4paper]{article}
\usepackage[utf8]{inputenc}
\usepackage[T1]{fontenc}
\usepackage{amsmath,amssymb}
\usepackage{amsthm}
\usepackage{hyperref}
\usepackage{geometry}
\geometry{margin=2.5cm}
\usepackage{enumitem}
\usepackage{booktabs}
\newtheorem{theorem}{Theorem}[section]
\newtheorem{lemma}[theorem]{Lemma}

\title{\bfseries Formalising the Trace–Determinant Link via the Stieltjes Transform\\
{\large A Research Programme for the GDST Proof of the Riemann Hypothesis}}
\author{Victor Geere \and Mentor}
\date{May 2026}

\begin{document}

\maketitle

\begin{abstract}
In the GDST (Greedy Dirichlet Sub‑Sum Transform) proof of the Riemann Hypothesis, 
the final step requires identifying the closed support of the spectral measure 
$\mu_\sigma$ with the set $\{\,t\in\mathbb{R}\mid\zeta(\tfrac12+it)=0\,\}$.  
This identification is achieved by expressing the Stieltjes transform of $\mu_\sigma$ 
in terms of the non‑trivial zeros of $\zeta(s)$ via a trace identity, and then 
applying analytic continuation. 

We outline a multi‑stage formalisation programme in Isabelle/HOL to carry out this 
argument.  The project involves building a theory of Stieltjes transforms of discrete 
measures, proving the asymptotic expansion relating moments to the transform, 
substituting the moment sums provided by the existing trace axioms, justifying the 
interchange of sums, demonstrating that the error functions sum to an entire function, 
and finally deriving the analytic continuation that forces all zeros onto the critical 
line.  The completion of this programme will turn the GDST proof from a conditional 
axiomatic proof into a fully verified theorem.
\end{abstract}

\section{Introduction}
The GDST programme~\cite{Geere2026} reduces the Riemann Hypothesis to the statement that 
the spectral measure $\mu_\sigma$ associated with the operator $T_\sigma$ has its 
closed support equal to the set of ordinates $t$ such that $\zeta(\tfrac12+it)=0$.  
The measure $\mu_\sigma$ is a purely discrete measure on $\mathbb{R}$, defined as 
$\sum_i \delta_{\lambda_{\sigma,i}}$ where $\lambda_{\sigma,i}$ are the eigenvalues of 
$T_\sigma$.  The available axioms (AX\_tr, AX\_zm) give the power sums 
$\sum_i \lambda_{\sigma,i}^k$ as
\[
\sum_i \lambda_{\sigma,i}^k = \sum_{\rho \in Z} \frac{1}{(\rho-\sigma)^k} 
+ E_k(\sigma) + G_k(\sigma),
\]
where $Z$ is the set of non‑trivial zeros of $\zeta(s)$ and $E_k, G_k$ are summable 
correction terms that combine to an entire function.  

The link to the zeros is made through the \emph{Stieltjes transform} of $\mu_\sigma$:
\[
S_\sigma(z) = \int_{\mathbb{R}} \frac{d\mu_\sigma(x)}{z - x}, \qquad z \notin \operatorname{supp}\mu_\sigma.
\]
For $|z|$ large, $S_\sigma(z)$ admits an asymptotic expansion in inverse powers whose 
coefficients are the moments $m_k(\sigma) = \sum_i \lambda_{\sigma,i}^k$.  Using the 
trace identities, the expansion becomes
\[
S_\sigma(z) \sim \sum_{k=1}^\infty \left( \sum_{\rho} \frac{1}{(\rho-\sigma)^k} \right) z^{-k-1}
= \sum_{\rho} \frac{1}{z - (\rho-\sigma)} + \text{entire}.
\]
By shifting the variable, we obtain a meromorphic function with poles exactly at the 
zeros $\rho$.  Analytic continuation and the fact that $\mu_\sigma$ is a real measure 
force these poles to lie on the real line, i.e., $\Im \rho = 0$.  This yields the 
desired conclusion that all non‑trivial zeros satisfy $\Re \rho = \frac12$.

This research programme provides a detailed plan for formalising the entire chain of 
reasoning inside Isabelle/HOL, from the definition of the Stieltjes transform to the 
final geometric conclusion.

\section{Objectives}
\begin{enumerate}[label=(\arabic*)]
    \item Develop a general theory of Stieltjes transforms of finite discrete measures 
          on $\mathbb{R}$.
    \item Prove the asymptotic expansion linking the Stieltjes transform to the 
          moment series in a half‑plane $\{|z| > R\}$.
    \item Instantiate the theory with the measure $\mu_\sigma$ and substitute the 
          moment sums using the trace axioms (AX\_tr, AX\_zm).
    \item Justify the interchange of infinite sums necessary to combine the moment 
          series with the sums over zeros.
    \item Show that the correction functions $E_k$ and $G_k$ sum to an entire function, 
          thereby demonstrating that the Stieltjes transform has the same singularities 
          as $\sum_\rho 1/(\rho - z)$.
    \item Prove that a discrete measure on $\mathbb{R}$ whose Stieltjes transform has 
          only real poles must itself be supported on the same set of real points; 
          apply this to $\mu_\sigma$ to deduce that all $\rho$ are real.
    \item Integrate the resulting theorem into the GDST theory, eliminating the 
          remaining gap in the final RH proof.
\end{enumerate}

\section{Mathematical Background}

\subsection{The Stieltjes transform of a discrete measure}
Let $\mu = \sum_{i} c_i \delta_{a_i}$ be a finite discrete measure on $\mathbb{R}$ with 
atoms at $\{a_i\}$ and weights $c_i>0$.  Its Stieltjes transform is
\[
S_\mu(z) = \int_{\mathbb{R}} \frac{d\mu(x)}{z - x} = \sum_i \frac{c_i}{z - a_i}, 
\qquad z \in \mathbb{C} \setminus \{a_i\}.
\]
This function is analytic on $\mathbb{C} \setminus \operatorname{supp}\mu$ and has 
simple poles at each atom $a_i$.

\subsection{Asymptotic expansion}
For $|z| > \sup_i |a_i|$, the geometric series expansion gives
\[
\frac{c_i}{z - a_i} = \sum_{k=0}^\infty \frac{c_i a_i^k}{z^{k+1}}.
\]
Summing over $i$ yields
\[
S_\mu(z) = \sum_{k=0}^\infty \frac{m_k}{z^{k+1}}, \qquad 
m_k = \int x^k d\mu = \sum_i c_i a_i^k,
\]
with uniform convergence on compact subsets of $\{|z| > R\}$ for any $R > \sup_i |a_i|$.  
This identity is the key bridge between moments and the analytic continuation of the 
Stieltjes transform.

\subsection{Connection to the GDST}
In our setting, $\mu = \mu_\sigma$ has $c_i = 1$ and $a_i = \lambda_{\sigma,i}$.  
Axioms AX\_tr and AX\_zm provide expressions for $m_k(\sigma)$ in terms of sums over 
the non‑trivial zeros $\rho$.  Specifically,
\[
m_k(\sigma) = \sum_{\rho} \frac{1}{(\rho-\sigma)^k} + H_k(\sigma),
\]
where $H_k(\sigma) = E_k(\sigma) + G_k(\sigma)$ are such that 
$\sum_{k=1}^\infty \frac{H_k(\sigma)}{z^{k+1}}$ converges to an entire function 
$\Phi_\sigma(z)$.  Consequently,
\[
S_{\mu_\sigma}(z) = \sum_{k=1}^\infty \frac{m_k(\sigma)}{z^{k+1}} 
= \sum_{\rho} \sum_{k=1}^\infty \frac{1}{(\rho-\sigma)^k z^{k+1}} + \Phi_\sigma(z).
\]
Exchanging the sums gives
\[
S_{\mu_\sigma}(z) = \sum_{\rho} \frac{1}{z - (\rho-\sigma)} + \Phi_\sigma(z).
\]
After the change of variable $w = z + \sigma$, we obtain a meromorphic function 
with poles exactly at $\rho$.

\subsection{Analytic continuation and the real‑pole argument}
Since $\mu_\sigma$ is a real measure, $S_{\mu_\sigma}(z)$ is real‑analytic on 
$\mathbb{C} \setminus \mathbb{R}$ and satisfies $\overline{S_{\mu_\sigma}(\bar z)} = 
S_{\mu_\sigma}(z)$.  Moreover, any pole of $S_{\mu_\sigma}$ must lie on the real axis 
because the only singularities of a discrete real measure's Stieltjes transform are 
the atoms themselves.  Consequently, the poles $\rho$ that appear in the meromorphic 
representation must be real.  Hence all non‑trivial zeros have $\Im \rho = 0$, i.e. 
$\Re \rho = \frac12$.

\section{Formalisation Plan}

\subsection{Phase 1 – Stieltjes transform of discrete measures}
\begin{itemize}
    \item Define a type for discrete measures on $\mathbb{R}$ with finite support 
          (or with summable weights), using a constructor that stores a countable set 
          of atoms and weights.
    \item Define the Stieltjes transform as a function $S: \mathbb{C} \to \mathbb{C}$,
          $\displaystyle S(z) = \sum_{a \in \text{atoms}} \frac{w(a)}{z - a}$, and 
          prove that it is holomorphic on $\mathbb{C} \setminus \overline{\{a_i\}}$.
    \item Prove elementary properties: real‑analyticity, invariance under complex 
          conjugation, and the fact that the poles are exactly the atoms.
\end{itemize}

\subsection{Phase 2 – Moment expansion}
\begin{itemize}
    \item For a discrete measure with atoms bounded in absolute value by $R$, prove 
          that for $|z| > R$,
          \[
          S(z) = \sum_{k=0}^\infty \frac{m_k}{z^{k+1}},
          \]
          where $m_k = \sum_i w_i a_i^k$ (the moments).  The proof uses the geometric 
          series and uniform convergence for $|z| \ge R' > R$.
    \item Establish absolute convergence and the ability to exchange the two sums 
          (over atoms and over powers) in the domain $|z| > R$.  This is a standard 
          Fubini/Tonelli application for sums.
\end{itemize}

\subsection{Phase 3 – Substitution of the GDST trace identities}
\begin{itemize}
    \item Instantiate the moment expansion for $\mu_\sigma$ (which satisfies the 
          necessary boundedness by the super‑exponential estimates from Part~I).
    \item Use AX\_tr and AX\_zm to rewrite $m_k(\sigma)$ as $\sum_\rho (\rho-\sigma)^{-k} 
          + H_k(\sigma)$.
    \item This yields
          \[
          S_{\mu_\sigma}(z) = \sum_{k=1}^\infty \frac{1}{z^{k+1}} \sum_{\rho} 
          \frac{1}{(\rho-\sigma)^k} + \sum_{k=1}^\infty \frac{H_k(\sigma)}{z^{k+1}}.
          \]
\end{itemize}

\subsection{Phase 4 – Interchange of sums over zeros and moments}
\begin{itemize}
    \item The double sum $\sum_{k=1}^\infty \sum_{\rho} \frac{1}{(\rho-\sigma)^k z^{k+1}}$ 
          must be commuted.  This is the most delicate point because the sum over 
          zeros is only conditionally convergent.  We require the assumption that the 
          series $\sum_\rho \frac{1}{|\rho-\sigma|}$ converges conditionally (which is 
          true for the non‑trivial zeros, but far from trivial).  In the context of 
          the GDST, we may rely on classical estimates from analytic number theory 
          (e.g., the density lemma $N(T) \sim \frac{T}{2\pi}\log T$) that guarantee 
          the summability of $|\rho-\sigma|^{-2}$ for $\sigma>1/2$, allowing the 
          interchange via Fubini's theorem after taking absolute values inside an 
          appropriate region.
    \item Prove that under the existing axiomatic framework (or an additional mild 
          absolute‑convergence axiom for the zero‑sum) the interchange is justified, 
          leading to
          \[
          \sum_{\rho} \sum_{k=1}^\infty \frac{1}{(\rho-\sigma)^k z^{k+1}} 
          = \sum_{\rho} \frac{1}{z - (\rho-\sigma)}.
          \]
\end{itemize}

\subsection{Phase 5 – Entirety of the error function}
\begin{itemize}
    \item The functions $H_k(\sigma)$ satisfy $\sum_{k} |H_k(\sigma)| < \infty$ (by 
          AX\_tr and AX\_zm).  Prove that this implies the series 
          $\Phi_\sigma(z) = \sum_{k=1}^\infty \frac{H_k(\sigma)}{z^{k+1}}$ converges 
          uniformly on compact sets to an entire function (since each term is analytic 
          and the tail is uniformly small for large $|z|$).
    \item Thus $S_{\mu_\sigma}(z) = \sum_{\rho} \frac{1}{z - (\rho-\sigma)} + \Phi_\sigma(z)$,
          with $\Phi_\sigma$ entire.
\end{itemize}

\subsection{Phase 6 – Analytic continuation and conclusion}
\begin{itemize}
    \item Show that $S_{\mu_\sigma}(z)$ is meromorphic on $\mathbb{C}$, with poles 
          exactly at the points $\rho-\sigma$, because $\Phi_\sigma$ is entire and the 
          sum $\sum_\rho \frac{1}{z - (\rho-\sigma)}$ is a meromorphic function with 
          those poles (this requires a standard Mittag‑Leffler argument for the zeros 
          of $\zeta$, which we may assume via another axiom or a classical theorem 
          formalised elsewhere).
    \item Prove that the Stieltjes transform of a real discrete measure has only real 
          poles (the atoms).  Therefore all points $\rho-\sigma$ must be real, i.e. 
          $\rho \in \mathbb{R}$.  Since the non‑trivial zeros are known to lie in the 
          critical strip $0<\Re s<1$, this forces $\Re \rho = 1/2$.
    \item Integrate this result into the GDST theory to replace the final `sorry` 
          and complete the proof of the Riemann Hypothesis.
\end{itemize}

\section{Resources and Existing Theories}
The formalisation will build on the following Isabelle libraries:
\begin{itemize}
    \item \texttt{HOL-Analysis}: measures, Lebesgue/Bochner integral, infinite sums, 
          complex analysis (holomorphic functions, analytic continuation).
    \item \texttt{Zeta\_Function}: definition of $\zeta(s)$ and its basic properties, 
          including the functional equation.
    \item \texttt{Countable} and \texttt{Set\_Integral}: discrete sums and integrals.
    \item The GDST theory file itself, which provides the spectral measure $\mu_\sigma$,
          the trace axioms, and the Fredholm determinant identity.
\end{itemize}
No additional AFP entries are strictly necessary, though the \texttt{Prime\_Number\_Theorem} 
entry could supply estimates for the zero density if we choose to avoid additional axioms.

\section{Timeline and Effort}
The following is a realistic estimate for a single researcher working part‑time (50\% FTE) 
with moderate experience in Isabelle/HOL analysis.

\begin{center}
\begin{tabular}{@{}ll@{}}
\toprule
\textbf{Phase} & \textbf{Effort} \\
\midrule
1. Stieltjes transform basics & 3–4 weeks \\
2. Moment expansion & 2–3 weeks \\
3. Substitution of trace identities & 1–2 weeks \\
4. Interchange of sums over zeros & 4–6 weeks (hard) \\
5. Entirety of error function & 2–3 weeks \\
6. Analytic continuation and final deduction & 3–4 weeks \\
\midrule
\textbf{Total} & $\approx$ 15–22 weeks \\
\bottomrule
\end{tabular}
\end{center}

The interchange of sums (Phase~4) is the most mathematically challenging part; 
it may require incorporating external estimates for the Riemann zeta zeros or adding 
a carefully justified axiom.  Tight collaboration with a number theorist is recommended 
during this phase.

\section{Conclusion}
This research programme lays out the path to formalise the trace–determinant link, the 
crucial missing piece in the GDST proof of the Riemann Hypothesis.  The formalisation is 
ambitious but well‑defined; once completed, it will turn the current axiomatic conditional 
proof into a self‑contained, fully verified theorem.  The programme also produces reusable 
libraries for Stieltjes transforms and moment problems that will benefit other formal 
analysis projects.

\begin{thebibliography}{9}
\bibitem{Geere2026}
V.~Geere, \emph{GDST Programme: Proof of the Riemann Hypothesis via the Greedy Dirichlet 
Sub‑Sum Transform}, Isabelle theory file, 2026.
\bibitem{Simon2011}
B.~Simon, \emph{Szegő's Theorem and its Descendants}, Princeton Univ. Press, 2011.
\bibitem{Hadamard}
J.~Hadamard, \emph{Sur la distribution des zéros de la fonction $\zeta(s)$ et ses 
conséquences arithmétiques}, Bull. Soc. Math. France, 1896.
\bibitem{Apostol}
T.~M.~Apostol, \emph{Introduction to Analytic Number Theory}, Springer, 1976.
\end{thebibliography}

\end{document}